% GENERATED FILE. Edit canonical lessons, then run npm run generate:books.
% canonical-source-hash: fnv1a64:e4ad8faa45a2ea10
% canonical-lessons: GE-C442-glatteis, GE-C442-raeumen, GE-C442-streuen, GE-C442-schneepflug, GE-C442-rutschen

\chapter{Schnee in der Stadt}
\label{ch:ge-schnee-in-der-stadt}
\hlchaptermodality{442}

\begin{chapteropening}
\textbf{By the end of this chapter:} \emph{I can talk about a town in winter: black ice, clearing snow, gritting, the snowplough and slipping.}

Complete the last lesson of chapter 442: I can talk about a town in winter: black ice, clearing snow, gritting, the snowplough and slipping.
\end{chapteropening}

\section[das Glatteis]{das Glatteis — the black ice, the icy road}

\label{lesson:GE-C442-glatteis}



\begin{quote}
Before the new one: say the German for the covered market, then the German for the pillar.
\end{quote}

\subsection*{You'll want to know: das Glatteis}
\textbf{das Glatteis} — \textquotedblleft{}the black ice, the icy road\textquotedblright{}.

\begin{quote}
Wegen Glatteis fahren heute keine Busse.
\end{quote}

\begin{quote}
\emph{There are no buses today because of the black ice.}
\end{quote}

No plural. \emph{Jemanden aufs Glatteis führen} — to trick someone.

\begin{cousinweb}
\emph{Glatteis} is built from \emph{glatt} 'smooth', inherited from Proto-Germanic \emph{gladaz}, and \emph{Eis} 'ice', from Proto-Germanic \emph{īsą} (EtymDB).
\end{cousinweb}

\begin{grammarlens}[title={What you've built}]
One more word for telling, explaining and describing at B1.
\end{grammarlens}

\subsection*{Your turn}
\begin{itemize}
\raggedright
  \item \emph{Say it:} \emph{das Glatteis}
  \item \emph{Say it:} \emph{das Glatteis}, once more
  \item \emph{Say it:} a sentence of your own with \emph{das Glatteis}
  \item \emph{Recall:} say the German for the covered market, then the German for the pillar, then say \emph{das Glatteis} again
\end{itemize}

\begin{tcolorbox}[breakable,colback=teal!4,colframe=teal!35!black,title={Before you move on}]
What does das Glatteis mean? (\textquotedblleft{}The black ice, the icy road\textquotedblright{}.) Say it once more.
\end{tcolorbox}
\section[räumen]{räumen — to clear (snow, a room)}

\label{lesson:GE-C442-raeumen}



\begin{quote}
Before the new one: say the German for the pillar, then the German for the black ice.
\end{quote}

\subsection*{You'll want to know: räumen}
\textbf{räumen} — \textquotedblleft{}to clear (snow, a room)\textquotedblright{}.

\begin{quote}
Am Morgen räumt der Nachbar den Schnee vor dem Haus.
\end{quote}

\begin{quote}
\emph{In the morning the neighbour clears the snow in front of the house.}
\end{quote}

\emph{Aufräumen} — to tidy up.

\begin{cousinweb}
\emph{räumen} is inherited from Old High German \emph{rūmen} 'to make room' and Proto-Germanic \emph{rūmijaną}, built on \emph{rūmaz} 'roomy' (EtymDB). German \emph{Raum} and English \emph{room} are cousins.
\end{cousinweb}

\begin{grammarlens}[title={What you've built}]
One more word for telling, explaining and describing at B1.
\end{grammarlens}

\subsection*{Your turn}
\begin{itemize}
\raggedright
  \item \emph{Say it:} \emph{räumen}
  \item \emph{Say it:} \emph{räumen}, once more
  \item \emph{Say it:} a sentence of your own with \emph{räumen}
  \item \emph{Recall:} say the German for the pillar, then the German for the black ice, then say \emph{räumen} again
\end{itemize}

\begin{tcolorbox}[breakable,colback=teal!4,colframe=teal!35!black,title={Before you move on}]
What does räumen mean? (\textquotedblleft{}To clear (snow, a room)\textquotedblright{}.) Say it once more.
\end{tcolorbox}
\section[streuen]{streuen — to grit (a road); to scatter, to sprinkle}

\label{lesson:GE-C442-streuen}



\begin{quote}
Before the new one: say the German for the black ice, then the German for to clear.
\end{quote}

\subsection*{You'll want to know: streuen}
\textbf{streuen} — \textquotedblleft{}to grit (a road); to scatter, to sprinkle\textquotedblright{}.

\begin{quote}
Die Stadt streut Salz auf die Straßen.
\end{quote}

\begin{quote}
\emph{The town spreads salt on the roads.}
\end{quote}

\emph{Gestreut} — gritted. \emph{Zucker auf den Kuchen streuen} — to sprinkle sugar on the cake.

\begin{cousinweb}
\emph{streuen} is inherited from Middle High German \emph{ströuwen} (EtymDB), and English \emph{strew} is a cousin.
\end{cousinweb}

\begin{grammarlens}[title={What you've built}]
One more word for telling, explaining and describing at B1.
\end{grammarlens}

\subsection*{Your turn}
\begin{itemize}
\raggedright
  \item \emph{Say it:} \emph{streuen}
  \item \emph{Say it:} \emph{streuen}, once more
  \item \emph{Say it:} a sentence of your own with \emph{streuen}
  \item \emph{Recall:} say the German for the black ice, then the German for to clear, then say \emph{streuen} again
\end{itemize}

\begin{tcolorbox}[breakable,colback=teal!4,colframe=teal!35!black,title={Before you move on}]
What does streuen mean? (\textquotedblleft{}To grit (a road); to scatter, to sprinkle\textquotedblright{}.) Say it once more.
\end{tcolorbox}
\section[der Schneepflug]{der Schneepflug — the snowplough (US: snowplow)}

\label{lesson:GE-C442-schneepflug}



\begin{quote}
Before the new one: say the German for to clear, then the German for to grit.
\end{quote}

\subsection*{You'll want to know: der Schneepflug}
\textbf{der Schneepflug} — \textquotedblleft{}the snowplough (US: snowplow)\textquotedblright{}.

\begin{quote}
In der Nacht fährt der Schneepflug durch unsere Straße.
\end{quote}

\begin{quote}
\emph{In the night the snowplough goes along our street.}
\end{quote}

Plural: \emph{die Schneepflüge}.

\begin{cousinweb}
\emph{Schneepflug} is built from \emph{Schnee} 'snow' and \emph{Pflug} 'plough', both inherited: \emph{Pflug} comes from Middle High German \emph{pfluoc}, a cousin of English \emph{plough} (EtymDB).
\end{cousinweb}

\begin{grammarlens}[title={What you've built}]
One more word for telling, explaining and describing at B1.
\end{grammarlens}

\subsection*{Your turn}
\begin{itemize}
\raggedright
  \item \emph{Say it:} \emph{der Schneepflug}
  \item \emph{Say it:} \emph{der Schneepflug}, once more
  \item \emph{Say it:} a sentence of your own with \emph{der Schneepflug}
  \item \emph{Recall:} say the German for to clear, then the German for to grit, then say \emph{der Schneepflug} again
\end{itemize}

\begin{tcolorbox}[breakable,colback=teal!4,colframe=teal!35!black,title={Before you move on}]
What does der Schneepflug mean? (\textquotedblleft{}The snowplough (US: snowplow)\textquotedblright{}.) Say it once more.
\end{tcolorbox}
\section[rutschen]{rutschen — to slip, to slide}

\label{lesson:GE-C442-rutschen}



\begin{quote}
Before the new one: say the German for to grit, then the German for the snowplough.
\end{quote}

\subsection*{You'll want to know: rutschen}
\textbf{rutschen} — \textquotedblleft{}to slip, to slide\textquotedblright{}.

\begin{quote}
Auf dem Eis rutscht sie und fällt hin.
\end{quote}

\begin{quote}
\emph{She slips on the ice and falls over.}
\end{quote}

Perfekt with \emph{sein}: \emph{ich bin gerutscht}. \emph{Die Rutsche} — the slide in a playground.

\begin{cousinweb}
\emph{rutschen} is an old expressive verb for sliding (EtymDB lists 'to slide'); its origin is not settled.
\end{cousinweb}

\begin{grammarlens}[title={What you've built}]
One more word for telling, explaining and describing at B1.
\end{grammarlens}

\subsection*{Your turn}
\begin{itemize}
\raggedright
  \item \emph{Say it:} \emph{rutschen}
  \item \emph{Say it:} \emph{rutschen}, once more
  \item \emph{Say it:} a sentence of your own with \emph{rutschen}
  \item \emph{Recall:} say the German for to grit, then the German for the snowplough, then say \emph{rutschen} again
\end{itemize}

\begin{tcolorbox}[breakable,colback=teal!4,colframe=teal!35!black,title={Before you move on}]
What does rutschen mean? (\textquotedblleft{}To slip, to slide\textquotedblright{}.) Say it once more.
\end{tcolorbox}
